% Appendix H: Data Flow Diagrams (DFD) Level 0, 1, 2
\chapter{DATA FLOW DIAGRAMS (DFD)}

This appendix provides a detailed narrative and structural description of the data flow across the Nebula platform, from the initial user request to the final cloud state persistence.

\section{DFD Level 0: Global Context Diagram}
1) \textbf{Entities}: User Browser, Nebula Control Plane, Cloud Provider APIs (AWS/Azure/GCP).
2) \textbf{Primary Flow}: The user sends "Infrastructure Intent" to the Control Plane. The Control Plane returns "Consolidated Inventory" and "Execution Logs."
3) \textbf{Cloud Interaction}: The Control Plane sends "Authenticated API Requests" to the providers and receives "Resource Metadata."

\section{DFD Level 1: Service Decomposition}
The Level 1 DFD breaks down the Control Plane into its core services: API Gateway, Task Broker, and worker Cluster.
1) \textbf{Flow 1 (Auth)}: User -> API (Verification) -> Database.
2) \textbf{Flow 2 (Provisioning)}: User -> API -> Redis (Task Queue) -> Worker (Terraform).
3) \textbf{Flow 3 (Logging)}: Worker -> Database (Log Storage) -> API (WebSocket) -> User UI.

\section{DFD Level 2: Task Runner Internal Flow}
This level focuses on the internal state transitions of the Celery Worker during a Terraform execution.
- \textbf{Stage A}: Receive Task ID and Vault Key.
- \textbf{Stage B}: Fetch Encrypted Credentials.
- \textbf{Stage C}: Generate HCL Template.
- \textbf{Stage D}: Initialize Terraform Backend.
- \textbf{Stage E}: Execute 'Apply' and capture STDOUT/STDERR.
- \textbf{Stage F}: Update Database and mark Task as 'Complete'.

\section{DFD Level 3: AI Copilot Analysis Flow}
The Level 3 DFD zooms into the failure-analysis sub-flow:
\begin{enumerate}
    \item Worker stores raw execution log after a failed task.
    \item API retrieves the failed task metadata and strips secrets or sensitive tokens.
    \item Sanitized log excerpt is combined with provider, region, and resource context.
    \item AI copilot receives the bounded prompt and returns a structured explanation.
    \item The explanation is stored and made visible to the user in the dashboard.
\end{enumerate}
This decomposition is important because it highlights the separation between raw operational data, sanitized analysis input, and user-facing remediation output.

\section{Data Stores and Trust Boundaries}
Across all DFD levels, the following trust boundaries should be recognized:
\begin{itemize}
    \item \textbf{Browser to API}: authenticated but untrusted client boundary.
    \item \textbf{API to Database}: trusted application boundary for persistent state.
    \item \textbf{API to Broker}: asynchronous handoff boundary.
    \item \textbf{Worker to Cloud APIs}: privileged execution boundary using decrypted provider credentials.
    \item \textbf{Worker to AI Analyzer}: sanitized diagnostic boundary without raw secrets.
\end{itemize}
Mapping these trust boundaries helps explain why the architecture isolates credential handling and why logs must be filtered before external analysis.

\newpage
% Appendix I: Detailed Test Reports and Bug Tracking
\chapter{DETAILED TEST REPORTS AND BUG TRACKING}

This appendix provides a comprehensive record of the testing phase, including 100+ test cases (TC-01 through TC-100) and a historical log of solved bugs.

\section{Exhaustive Test Case Matrix (TC-01 to TC-100)}
\begin{table}[H]
\centering
\caption{Nebula 100-Item Test Matrix (Summary)}
\label{tab:100_tests}
\begin{tabular}{|l|l|l|l|l|}
\hline
\textbf{Test ID} & \textbf{Category} & \textbf{Description} & \textbf{Pass/Fail} & \textbf{Notes} \\ \hline
TC-31 & UI & Dashboard Sidebar Collapse & PASS & - \\ \hline
TC-45 & Security & Brute Force Protection & PASS & 5 attempts limit \\ \hline
TC-52 & API & Invalid Resource Name Check & PASS & Returns 400 Bad Req \\ \hline
TC-68 & Sync & AWS Multi-Region Discovery & PASS & Verified us-east-1/2 \\ \hline
TC-75 & Vault & Key Rotation Logic & PASS & Atomic update \\ \hline
TC-88 & Network & VNet Peering Success & PASS & Azure backbone confirmed \\ \hline
TC-99 & Final & 24hr System Stability & PASS & 0 Memory leaks \\ \hline
\end{tabular}
\end{table}

\section{Test Execution Protocol}
The testing process followed a layered order so that failures could be localized quickly:
\begin{enumerate}
    \item validate schemas and input constraints at the unit level,
    \item test API routes with mocked database sessions,
    \item verify queue delivery and worker execution using integration tests,
    \item execute sandbox cloud tasks with low-cost resources,
    \item review dashboard behavior and log presentation from an operator perspective.
\end{enumerate}

\section{Representative Detailed Test Extract}
\begin{lstlisting}[language=SQL, caption=Detailed Test Extract]
TC-01 | Auth    | Login with valid credentials                  | PASS | JWT issued
TC-02 | Auth    | Login with invalid password                   | PASS | 401 returned
TC-03 | Vault   | Encrypt AWS access key                        | PASS | ciphertext stored
TC-04 | Vault   | Decrypt credential during worker execution    | PASS | plaintext not logged
TC-05 | API     | Reject malformed resource payload             | PASS | schema validation hit
TC-06 | Queue   | Enqueue provisioning task                     | PASS | broker accepted message
TC-07 | Queue   | Worker picks queued task                      | PASS | status moved to running
TC-08 | IaC     | Terraform init for AWS module                 | PASS | providers resolved
TC-09 | IaC     | Terraform apply for sample VM                 | PASS | resource active
TC-10 | Sync    | Refresh AWS inventory                         | PASS | canonical records updated
TC-11 | Sync    | Refresh Azure inventory                       | PASS | canonical records updated
TC-12 | Sync    | Refresh GCP inventory                         | PASS | canonical records updated
TC-13 | AI      | Detect AccessDenied signature                 | PASS | mapped to IAM category
TC-14 | AI      | Detect quota exhaustion                       | PASS | quota remediation shown
TC-15 | UI      | Show running status badge                     | PASS | badge transitions correct
TC-16 | UI      | Stop polling after task completion            | PASS | timer cleared
TC-17 | Audit   | Log delete action                             | PASS | activity row written
TC-18 | Recovery| Retry transient cloud timeout                 | PASS | next retry scheduled
\end{lstlisting}

\section{Historical Bug Log and Resolutions}
The following log tracks critical bugs identified during the integration testing phase and their subsequent fixes.

\begin{itemize}
    \item \textbf{Bug ID: B-102 (Credential Scrambling)}
        \begin{itemize}
            \item \textit{Issue}: Special characters in Azure secrets were causing Fernet decryption errors.
            \item \textit{Fix}: Implemented Base64 encoding/decoding before encryption/decryption to ensure character integrity.
        \end{itemize}
    \item \textbf{Bug ID: B-205 (Redis Queue Congestion)}
        \begin{itemize}
            \item \textit{Issue}: Bulk provisioning of 20+ VMs caused the Redis broker to timeout.
            \item \textit{Fix}: Increased the \texttt{visibility_timeout} in Celery and optimized the database connection pool in the worker cluster.
        \end{itemize}
    \item \textbf{Bug ID: B-309 (UI State Mismatch)}
        \begin{itemize}
            \item \textit{Issue}: Dashboard counts did not update automatically after a resource was deleted from the cloud portal directly.
            \item \textit{Fix}: Implemented a periodic "background reconcile" task that triggers every 5 minutes to keep the database in sync with actual cloud state.
        \end{itemize}
    \item \textbf{Bug ID: B-411 (Task Retry Storm)}
        \begin{itemize}
            \item \textit{Issue}: Retryable provider errors caused too many workers to retry in parallel, amplifying queue depth.
            \item \textit{Fix}: Added jitter to the exponential backoff calculation and reduced worker prefetch.
        \end{itemize}
    \item \textbf{Bug ID: B-507 (Unsafe Log Disclosure)}
        \begin{itemize}
            \item \textit{Issue}: Raw Terraform error output could include values that should not be forwarded to the AI analyzer.
            \item \textit{Fix}: Added sanitization rules for secrets, access keys, endpoint tokens, and private IP patterns before analysis.
        \end{itemize}
\end{itemize}

\newpage
